% =====================================================================
% 第十二讲 (Exposé XII) 代数几何与解析几何
% (Géométrie algébrique et géométrie analytique)
% M. Raynaud
% 原书页码: pp. 239-257
% =====================================================================
\chapter{代数几何与解析几何}

\noindent Mme M. Raynaud\footnote{按 A. Grothendieck 的未发表笔记.}

如 \cite{exp12-10} 那样进行，给复数域 $\mathbb{C}$ 上每个局部有限型概形 $X$ 关联一个解析空间 $X^{\mathrm{an}}$\index[notations]{$X^{\mathrm{an}}$}，它的底集是 $X(\mathbb{C})$.

在本讲的第 2 和第 3 号中，给出 $X$ 与 $X^{\mathrm{an}}$ 的通常性质之间、以及态射 $f \colon X \to Y$ 与关联态射 $f^{\mathrm{an}} \colon X^{\mathrm{an}} \to Y^{\mathrm{an}}$ 的性质之间的"字典".

然后证明 $X$ 与 $X^{\mathrm{an}}$ 上凝聚层之间的比较定理 (在 \cite[第 12 号]{exp12-10} 中对射影簇建立) 在 $X$ 是紧合概形时仍然成立.

最后在第 5 号中证明 $X$ 的有限平展覆盖范畴与 $X^{\mathrm{an}}$ 的有限平展覆盖范畴等价. 作为给读者的赠品，给出 \emph{Grauert-Remmert 定理} \cite{exp12-6} 的新证明，用\emph{奇点消解} \cite{exp12-8}.

\section{概形的关联解析空间}

设 $X$ 是 $\mathbb{C}$ 上局部有限型概形. 设 $\Phi$ 是从解析空间范畴 \cite[第 9 号]{exp12-4} 到集合范畴的函子，它把解析空间 $X$ 配给 $\mathbb{C}$-代数环化空间态射集 $\mathrm{Hom}_{\mathbb{C}}(X, X)$. 有下列定理：

\begin{theorem}[定理与定义]\label{thm:XII.1.1}
函子 $\Phi$ 被解析空间 $X^{\mathrm{an}}$ 和态射 $\phi \colon X^{\mathrm{an}} \to X$ 表示. 称 $X^{\mathrm{an}}$ 为 $X$ 的\emph{关联解析空间}\index[term]{guanlianjiexikongjian@关联解析空间}.
\end{theorem}

若 $|X^{\mathrm{an}}|$ 是 $X^{\mathrm{an}}$ 的底集，$\phi$ 诱导从 $|X^{\mathrm{an}}|$ 到 $X$ 的在 $\mathbb{C}$ 中取值的点集 $X(\mathbb{C})$ 的双射. 而且，对 $X^{\mathrm{an}}$ 的每个点 $x$，态射
\[
\phi_x \colon \mathcal{O}_{X,\phi(x)} \longrightarrow \mathcal{O}_{X^{\mathrm{an}},x},
\]
必然局部，经完备化给出同构
\[
\widehat{\phi_x} \colon \widehat{\mathcal{O}}_{X,\phi(x)} \xrightarrow{\sim} \widehat{\mathcal{O}}_{X^{\mathrm{an}},x}.
\]
特别态射 $\phi$ 平坦.

注意 $\phi$ 诱导 $X^{\mathrm{an}}$ 到 $X(\mathbb{C})$ 的双射的事实由 $X^{\mathrm{an}}$ 的万有性质得出. 另一方面有下列断言：

a) 若定理对概形 $Y$ 成立，则对 $Y$ 的任何子概形 $X$ 也成立. 先假设 $X$ 是 $Y$ 的开子概形；若 $\psi \colon Y^{\mathrm{an}} \to Y$ 是典范态射，$\psi^{-1}(X)$ 是 $Y^{\mathrm{an}}$ 的开集，配备 $Y^{\mathrm{an}}$ 诱导的解析空间结构. 由于解析空间 $X$ 到 $X$ 的每个态射由后者的万有性质经 $Y^{\mathrm{an}}$ 分解，从而经 $X^{\mathrm{an}}$ (即纤维积 $Y^{\mathrm{an}} \times_Y X$) 分解，$X^{\mathrm{an}}$ 是 $X$ 的关联解析空间. 最后关于 $\phi_x$ 的断言显然.

现在只需考虑 $X$ 是 $Y$ 的闭子概形的情形. 设 $I$ 是定义 $X$ 的 $\mathcal{O}_Y$-凝聚理想；则 $I \cdot \mathcal{O}_{Y^{\mathrm{an}}}$ 是 $\mathcal{O}_{Y^{\mathrm{an}}}$ 上的凝聚理想层，它定义 $Y^{\mathrm{an}}$ 的闭解析子空间 $X^{\mathrm{an}}$；如开子概形的情形那样看到 $X^{\mathrm{an}}$ 是 $X$ 的关联解析空间. 设 $\phi \colon X^{\mathrm{an}} \to X$ 是典范态射. 对 $X^{\mathrm{an}}$ 的每个点 $x$，态射 $\phi_x$ 正是由 $\psi_x$ 诱导的态射
\[
\mathcal{O}_{Y,\psi(x)}/I_{\psi(x)} \longrightarrow \mathcal{O}_{Y^{\mathrm{an}},x}/I_{\psi(x)} \cdot \mathcal{O}_{Y^{\mathrm{an}},x}
\]
的完备化
\[
\widehat{\phi_x} \colon \widehat{\mathcal{O}}_{Y,\psi(x)}/I_{\psi(x)} \cdot \widehat{\mathcal{O}}_{Y,\psi(x)} \longrightarrow \widehat{\mathcal{O}}_{Y^{\mathrm{an}},x}/I_{\psi(x)} \cdot \widehat{\mathcal{O}}_{Y^{\mathrm{an}},x}
\]
是同构，因为 $\widehat{\psi_x}$ 是同构，这证明 a).

b) 若有两个 $\mathbb{C}$-概形 $X_1$，$X_2$，使 $X_1^{\mathrm{an}}$ 和 $X_2^{\mathrm{an}}$ 存在，则 $(X_1 \times X_2)^{\mathrm{an}}$ 也存在. 设 $\phi_1 \colon X_1^{\mathrm{an}} \to X_1$，$\phi_2 \colon X_2^{\mathrm{an}} \to X_2$ 是典范态射，$p_1$，$p_2$ 是 $X_1^{\mathrm{an}} \times X_2^{\mathrm{an}}$ 的两个投影. 由 EGA I 1.8.1 形式地得出 $X_1 \times X_2$ 是 $X_1$ 和 $X_2$ 在局部环环化空间范畴中的积；由此得出态射 $\phi_1 \cdot p_1$ 和 $\phi_2 \cdot p_2$ 定义态射 $\phi \colon X_1^{\mathrm{an}} \times X_2^{\mathrm{an}} \to X_1 \times X_2$，且对 $(X_1^{\mathrm{an}} \times X_2^{\mathrm{an}}, \phi)$ 表示函子 $X \mapsto \mathrm{Hom}_{\mathbb{C}}(X, X_1 \times X_2)$.

c) 若 $E^1$ 表示维数 1 的仿射空间，i.e. 配备全纯函数\index[term]{quanchunhanshu@全纯函数}层的拓扑空间 $\mathbb{C}$，函子\footnote{[2003 年增补：$E^1_{\mathbb{C}}$ 表示 $\mathbb{C}$ 上的 (代数) 仿射直线.]} $X \mapsto \mathrm{Hom}_{\mathbb{C}}(X, E^1_{\mathbb{C}})$ 被 $E^1$ 表示，典范态射 $\phi \colon E^1 \to E^1_{\mathbb{C}}$ 是显然态射. 确实，给出解析空间 $X$ 到 $E^1_{\mathbb{C}}$ 的态射等价于给出 $\Gamma(X, \mathcal{O}_X)$ 的元素，后者也等价于给出 $X$ 到 $E^1$ 的态射. 显然有双射 $|E^1| \simeq E^1_{\mathbb{C}}(\mathbb{C})$，且对每个点 $x \in E^1$，态射 $\widehat{\phi_x}$ 正是 $\mathbb{C}$ 上一元形式幂级数环的恒等态射.

由 b) 和 c) 得出定理对仿射空间 $E^n_{\mathbb{C}}$ ($n \geqslant 0$) 成立. 用 a)，看到对 $\mathbb{C}$ 上每个局部有限型仿射概形 $X$ 也成立. 若不再假设 $X$ 仿射，且 $(X_i)$ 是 $X$ 的仿射开集覆盖，由万有性质和 a) 得出诸 $X_i^{\mathrm{an}}$ 粘合，从而定义 $X$ 的关联解析空间 $X^{\mathrm{an}}$.

\refstepcounter{theorem}\noindent\textbf{\thesection.\arabic{theorem}.}\quad 设 $f \colon X \to Y$ 是 $\mathbb{C}$ 上局部有限型概形的态射. 若 $\phi \colon X^{\mathrm{an}} \to X$ 和 $\psi \colon Y^{\mathrm{an}} \to Y$ 是典范态射，由 $Y^{\mathrm{an}}$ 的万有性质得出存在唯一态射 $f^{\mathrm{an}} \colon X^{\mathrm{an}} \to Y^{\mathrm{an}}$，使图式
\[
\begin{tikzcd}[row sep=large, column sep=large]
X^{\mathrm{an}} \arrow[r] \arrow[d,"{f^{\mathrm{an}}}"] & X \arrow[d,"{f}"] \\
Y^{\mathrm{an}} \arrow[r] & Y
\end{tikzcd}
\]
交换. 于是定义了从 $\mathbb{C}$ 上局部有限型概形范畴到解析空间范畴的函子 $\Phi$.

函子 $\Phi$ 与有限投射极限交换. 只需看到 $\Phi$ 与纤维积交换. 若 $X$，$Y$，$Z$ 是 $\mathbb{C}$ 上局部有限型概形，由 $X \times_Z Y$ 是 $X$ 和 $Y$ 在 $Z$ 上方的局部环环化空间范畴中的纤维积的事实得出 $X^{\mathrm{an}} \times_{Z^{\mathrm{an}}} Y^{\mathrm{an}}$ 满足刻画 $(X \times_Z Y)^{\mathrm{an}}$ 的万有性质.

\refstepcounter{theorem}\noindent\textbf{\thesection.\arabic{theorem}.}\quad 设 $X$ 是 $\mathbb{C}$ 上局部有限型概形，$X^{\mathrm{an}}$ 是关联解析空间，$\phi \colon X^{\mathrm{an}} \to X$ 是典范态射. 若 $F$ 是 $\mathcal{O}_X$-模，逆像 $\phi^{*} F = F^{\mathrm{an}}$ 是 $\mathcal{O}_{X^{\mathrm{an}}}$ 上的模层. 于是定义了从 $\mathcal{O}_X$-模范畴到 $X^{\mathrm{an}}$ 上模范畴的函子. 该函子与归纳极限交换 (EGA IV$_0$ 4.3.2). 层 $\mathcal{O}_{X^{\mathrm{an}}}$ 凝聚 \cite[第 18 号 §2 定理 2]{exp12-4}，故它把凝聚层变成\emph{凝聚层} (EGA IV$_0$ 5.3.11). 还有：

\begingroup
\addtocounter{theorem}{-1}
\renewcommand{\thetheorem}{\thesection.3.1}
\begin{proposition}\label{prop:XII.1.3.1}
把 $\mathcal{O}_{X^{\mathrm{an}}}$-模 $F$ 配给它在 $X^{\mathrm{an}}$ 上的逆像 $F^{\mathrm{an}}$ 的函子正合、忠实、保守.
\end{proposition}
\endgroup

正合性由态射 $\phi \colon X^{\mathrm{an}} \to X$ 平坦 (\cref{thm:XII.1.1}) 得出. 证明函子 $F \mapsto F^{\mathrm{an}}$ 忠实. 考虑到正合性，只需证明 $F^{\mathrm{an}}$ 为零时 $F$ 为零. 对 $X^{\mathrm{an}}$ 的每个点 $x$，有 $F_{\phi(x)} \otimes_{\mathcal{O}_{X,\phi(x)}} \mathcal{O}_{X^{\mathrm{an}},x} = 0$. 态射 $\mathcal{O}_{X,\phi(x)} \to \mathcal{O}_{X^{\mathrm{an}},x}$ 忠实平坦，故对 $X$ 的每个闭点 $\phi(x)$ 有 $F_{\phi(x)} = 0$，且 $X$ 是 Jacobson 概形 (EGA IV 10.4.8)，这蕴含 $F$ 为零.

函子 $F \mapsto F^{\mathrm{an}}$ 保守的事实由正合性和忠实性形式地得出.

\section{概形与关联解析空间的性质比较}

\begin{proposition}\label{prop:XII.2.1}
设 $X$ 是 $\mathbb{C}$ 上局部有限型概形，$X^{\mathrm{an}}$ 是关联解析空间，$n$ 是整数. 考虑性质 $P$：
\begin{enumerate}[i)]
\item 非空
\item[i')] 离散
\item[ii)] Cohen-Macaulay
\item[iii)] $(S_n)$
\item[iv)] 正则
\item[v)] $(R_n)$
\item[vi)] 正规
\item[vii)] 既约
\item[viii)] 维数 $n$.
\end{enumerate}
则 $X$ 具有性质 $P$ 当且仅当 $X^{\mathrm{an}}$ 具有性质 $P$.
\end{proposition}

设 $\phi \colon X^{\mathrm{an}} \to X$ 是典范态射. (i) 由 $|X^{\mathrm{an}}| = X(\mathbb{C})$ (\cref{thm:XII.1.1}) 和 $X$ 是 Jacobson 概形 (EGA IV 10.4.8) 得出. 说 $X$ (resp. $X^{\mathrm{an}}$) 离散等价于说 $\dim X = 0$ (resp. 由 \cite[第 19 号 §4 推论 6]{exp12-4}，$\dim X^{\mathrm{an}} = 0$)；故 (i') 由 (viii) 得出.

设 $P$ 是 (ii) 至 (vii) 之一. 为使 $X$ 具有性质 $P$，必要且充分条件是 $P$ 在 $X$ 的每个闭点成立；确实，$X$ 是 excellent (EGA IV 7.8.6 (iii))，故 $X$ 验证 $P$ 的点集开 (loc. cit.)，且该开集含所有闭点时等于整个 $X$. 说 $X$ (resp. $X^{\mathrm{an}}$) 具有性质 $P$ 于是等价于：对 $X^{\mathrm{an}}$ 的每个点 $x$，局部环 $\mathcal{O}_{X,\phi(x)}$ (resp. $\mathcal{O}_{X^{\mathrm{an}},x}$) 具有性质 $P$. 由于 excellent 局部环具有性质 $P$ 的事实可在完备化后看到，命题由同构 $\widehat{\mathcal{O}}_{X,\phi(x)} \xrightarrow{\sim} \widehat{\mathcal{O}}_{X^{\mathrm{an}},x}$ 在情形 (ii) 至 (vii) 中得出. 情形 (viii) 也如此，考虑到关系
\[
\dim X = \sup_x \dim \mathcal{O}_{X,\phi(x)}, \quad \dim X^{\mathrm{an}} = \sup_x \dim \mathcal{O}_{X^{\mathrm{an}},x},
\]
其中 $x \in X^{\mathrm{an}}$. 这完成证明.

\begin{proposition}\label{prop:XII.2.2}
设 $X$ 是 $\mathbb{C}$ 上局部有限型概形，$\phi \colon X^{\mathrm{an}} \to X$ 是典范态射，$T$ 是 $X$ 的局部可构部分. 则有关系
\[
\overline{\phi^{-1}(T)} = \phi^{-1}(\overline{T}).
\]
\end{proposition}

可假设 $T$ 是 $X$ 的稠密开集. 设 $H$ 是 $X$ 的底集为 $X - T$ 的既约闭子概形；关联空间 $H^{\mathrm{an}}$ 是 $X^{\mathrm{an}}$ 的底集为 $X^{\mathrm{an}} - \phi^{-1}(T)$ 的闭解析子空间. 需证明 $H^{\mathrm{an}}$ 的每个点 $x$ 属于 $\phi^{-1}(\overline{T})$. 在这样的点 $x$，解析空间芽 $(X^{\mathrm{an}}, x)$ 含子芽 $(H^{\mathrm{an}}, x)$，后者由 $\mathcal{O}_{X^{\mathrm{an}},x}$ 的非幂零理想定义. 于是由 Nullstellensatz \cite[第 19 号 §4 推论 3]{exp12-4} 得出 $x$ 的每个开邻域含有 $X^{\mathrm{an}}$ 的不属于 $H^{\mathrm{an}}$ 的点，这确实证明 $x \in \phi^{-1}(\overline{T})$.

\begin{corollary}\label{cor:XII.2.3}
设 $X$ 是 $\mathbb{C}$ 上局部有限型概形，$\phi \colon X^{\mathrm{an}} \to X$ 是典范态射，$T$ 是 $X$ 的局部可构部分. 为使 $T$ 是开部分 (resp. 闭部分，resp. 稠密部分)，必要且充分条件是 $\phi^{-1}(T)$ 如此.
\end{corollary}

该推论由 \ref{prop:XII.2.2} 和 $X$ 是 Jacobson 概形 (EGA IV 10.4.8)、故 $X$ 的两个在非常稠密集 $X(\mathbb{C})$ 上有同样迹的局部可构部分相等的事实得出.

\begin{proposition}\label{prop:XII.2.4}
设 $X$ 是 $\mathbb{C}$ 上局部有限型概形. 为使 $X$ 连通 (resp. 不可约)，必要且充分条件是 $X^{\mathrm{an}}$ 如此.
\end{proposition}

假设 $X^{\mathrm{an}}$ 连通 (resp. 不可约). $X^{\mathrm{an}}$ 在 $X$ 中的像 $X(\mathbb{C})$ 此时连通 (resp. 不可约). 由此得出 $X$ 连通 (resp. 不可约)，因为 $X$ 与 $X(\mathbb{C})$ 的闭部分一一对应 (EGA IV 10.1.2).

反之假设 $X$ 连通 (resp. 不可约)，证明 $X^{\mathrm{an}}$ 亦然. 可限于 $X$ 不可约的情形. 假设 $X$ 连通. 给定 $X$ 的点 $x$，存在 $X$ 的不可约闭子概形 $X_1, \ldots, X_n$ 的有限序列 (其中 $x \in X_1$，$y \in X_n$，且对 $1 \leqslant i \leqslant n - 1$ 有 $X_i \cap X_{i+1} \neq \varnothing$) 的点 $y \in X$ 的集同时开且闭，故等于整个 $X$. 对前述那样的序列 $X_1, \ldots, X_n$，对 $1 \leqslant i \leqslant n - 1$ 也有 $X_i^{\mathrm{an}} \cap X_{i+1}^{\mathrm{an}} \neq \varnothing$；若已证诸 $X_i^{\mathrm{an}}$ 连通，则 $X^{\mathrm{an}}$ 也连通.

此后假设 $X$ 不可约. 还可假设 $X$ 仿射. 若 $(U_i)_{i \in I}$ 是 $X$ 的仿射开集覆盖，其中任意两个开集的交非空，则同样的性质对 $X^{\mathrm{an}}$ 的覆盖 $(U_i^{\mathrm{an}})_{i \in I}$ 成立；若已证诸 $U_i^{\mathrm{an}}$ 不可约，则 $X^{\mathrm{an}}$ 也如此.

还可假设 $X$ 正规. 设 $\widetilde{X}$ 是 $X$ 的正规化；由于态射 $\widetilde{X} \to X$ 满射，$\widetilde{X}^{\mathrm{an}} \to X^{\mathrm{an}}$ 也满射，这证明 $\widetilde{X}^{\mathrm{an}}$ 不可约时 $X^{\mathrm{an}}$ 也如此.

此后假设 $X$ 仿射正规. 由于 $X^{\mathrm{an}}$ 的局部环整，说 $X^{\mathrm{an}}$ 不可约与说它连通相同. 确实，若 $F$ 是 $X^{\mathrm{an}}$ 的闭解析部分，其中 $\mathrm{codim}_x(F, X^{\mathrm{an}}) = 0$ 的点 $x$ 的集是 $X^{\mathrm{an}}$ 的闭解析子集 \cite[第 20 号 A 推论 1]{exp12-4}，也开；若 $X^{\mathrm{an}}$ 连通，这证明 $F \neq X^{\mathrm{an}}$ 时 $F$ 稀疏，故 $X^{\mathrm{an}}$ 不可约. 于是归结到证明 $X^{\mathrm{an}}$ 连通.

设
\[
i \colon X \longrightarrow P
\]
是 $X$ 的紧化，其中 $P$ 是射影正规 $\mathbb{C}$-概形，$i$ 是支配开浸入. 由 \cite[第 12 号 定理 1]{exp12-10} 得出 $P^{\mathrm{an}}$ 连通. 由于 $X^{\mathrm{an}}$ 是从 $P^{\mathrm{an}}$ 去掉稀疏闭解析部分得到的，由下面 \ref{lem:XII.2.5} 得出 $X^{\mathrm{an}}$ 也连通.

\begin{lemma}\label{lem:XII.2.5}
设 $P$ 是连通正规解析空间，$Y$ 是稀疏闭解析部分，则 $X = P - Y$ 连通.
\end{lemma}

当 $Y$ 余维数 $\geqslant 2$ 时，命题由 \cite[第 3 号 命题 4]{exp12-11} 得出. 在一般情形，可以假设 (把 $P$ 中余维数 $\geqslant 2$ 的闭解析部分去掉) $P$ 和 $Y$ (看作 $P$ 的既约解析子空间) 正则. 由隐函数定理，$Y$ 的每个点 $y$ 有同构于仿射空间 $E^n$ 的球的邻域 $U$，使 $U \cap Y$ 由若干坐标函数的消没定义. 这证明 $U - U \cap Y$ 连通，故 $X$ 也连通.

\begin{corollary}\label{cor:XII.2.6}
设 $X$ 是 $\mathbb{C}$ 上局部有限型概形；由典范态射 $X^{\mathrm{an}} \to X$ 诱导的态射
\[
\pi_0(X^{\mathrm{an}}) \longrightarrow \pi_0(X)
\]
双射.
\end{corollary}

\section{态射性质的比较}

\begin{proposition}\label{prop:XII.3.1}
设 $f \colon X \to Y$ 是 $\mathbb{C}$ 上局部有限型概形的态射，$f^{\mathrm{an}} \colon X^{\mathrm{an}} \to Y^{\mathrm{an}}$ 是 $f$ 在关联解析空间上导出的态射. 设 $P$ 是下列性质：
\begin{enumerate}[i)]
\item 平坦
\item[ii)] 非分歧 (i.e. 非分歧)
\item[iii)] 平展
\item[iv)] 光滑
\item[v)] 正规
\item[vi)] 既约
\item[vii)] 单射
\item[viii)] 分离
\item[ix)] 同构
\item[x)] 单态射
\item[xi)] 开浸入.
\end{enumerate}
则 $f$ 具有性质 $P$ 当且仅当 $f^{\mathrm{an}}$\index[notations]{$f^{\mathrm{an}}$} 具有性质 $P$.
\end{proposition}

记 $\phi \colon X^{\mathrm{an}} \to X$ 和 $\psi \colon Y^{\mathrm{an}} \to Y$ 为典范态射. 设 $x$ 是 $X^{\mathrm{an}}$ 的点，$y = f^{\mathrm{an}}(x)$. 由 $f$ 和 $f^{\mathrm{an}}$ 导出的态射 $\mathcal{O}_{Y^{\mathrm{an}},y} \to \mathcal{O}_{X^{\mathrm{an}},x}$ 和 $\mathcal{O}_{Y,\psi(y)} \to \mathcal{O}_{X,\phi(x)}$ 经完备化给出同样态射 (\cref{thm:XII.1.1}). 由 \cite[第 3 章 §5 命题 4]{exp12-2} (resp. EGA IV 17.4.4)，说 $f^{\mathrm{an}}$ 验证性质 (i) (resp. (ii)) 与说 $f$ 在 $X$ 的每个闭点验证 (i) (resp. (ii)) 相同. 由于 $X$ 中验证 (i) (resp. (ii)) 的点集开 (EGA IV 11.1.1 和 I 3.3)，这证明 (i) 和 (ii)，从而 (iii).

设 $P$ 是性质 (iv) (resp. (v)，resp. (vi)). 考虑到 \ref{prop:XII.2.1} ((v)，(vi)，(vii))，说 $f^{\mathrm{an}}$ 在 $Y^{\mathrm{an}}$ 的各点 $y$ 处的几何纤维正则 (resp. 正规，resp. 既约) 与说 $f$ 在 $Y$ 的各闭点 $\psi(y)$ 处的几何纤维如此相同. 情形 (iv) (resp. (v)，resp. (vi)) 于是由 (i) 和 $Y$ 中 $f$ 的几何纤维正则的点集开 (EGA IV 12.1.7) 的事实得出.

(vii). 若 $f$ 单射，$f^{\mathrm{an}}$ 也单射. 反之假设 $f^{\mathrm{an}}$ 单射，证明 $f$ 也单射. 可假设 $f$ 有限型. 由于 $f^{\mathrm{an}}$ 单射，$f$ 在 $Y$ 的闭点处的纤维根式；由于 $Y$ 中纤维根式的点集局部可构 (EGA IV 9.6.1)，且 $Y$ 是 Jacobson 概形，$f$ 的所有纤维根式故单射.

(viii). 设 $\Delta \colon X \to X \times_Y X$ 和 $\Delta^{\mathrm{an}} \colon X^{\mathrm{an}} \to X^{\mathrm{an}} \times_{Y^{\mathrm{an}}} X^{\mathrm{an}}$ 是对角浸入，$\Theta \colon X^{\mathrm{an}} \times_{Y^{\mathrm{an}}} X^{\mathrm{an}} \to X \times_Y X$ 是典范态射. 由 \ref{cor:XII.2.3}，说 $\Delta(X)$ 在 $X \times_Y X$ 中闭与说 $\Delta^{\mathrm{an}}(X^{\mathrm{an}})$ 在 $X^{\mathrm{an}} \times_{Y^{\mathrm{an}}} X^{\mathrm{an}}$ 中闭相同.

由于开浸入正是平展单射态射 (EGA IV 17.9.1 和 \cite[第 13 号 §1]{exp12-4})，(xi) 由 (iii) 和 (vii) 得出. 同构正是满射开浸入，故 (ix) 由 (xi) 和下面 (3.2 (i)) 得出. 说 $f$ 是单态射等价于说对角态射 $\Delta \colon X \to X \times_Y X$ 是同构，故 (x) 由 (ix) 得出.

\begin{proposition}\label{prop:XII.3.2}
设 $X$ 和 $Y$ 是 $\mathbb{C}$ 上两个局部有限型概形，$f \colon X \to Y$ 是有限型态射，$f^{\mathrm{an}} \colon X^{\mathrm{an}} \to Y^{\mathrm{an}}$ 是 $f$ 在关联解析空间上导出的态射. 设 $P$ 是下列性质：
\begin{enumerate}[i)]
\item 满射
\item[ii)] 支配
\item[iii)] 闭浸入
\item[iv)] 浸入
\item[v)] 紧合\footnote{称解析空间态射紧合，若按 \cite[第 1 章 §10 第 1 号]{exp12-1} 的意义紧合且分离.}
\item[vi)] 有限.
\end{enumerate}
则 $f$ 具有性质 $P$ 当且仅当 $f^{\mathrm{an}}$ 具有性质 $P$.
\end{proposition}

设 $\phi \colon X^{\mathrm{an}} \to X$ 和 $\psi \colon Y^{\mathrm{an}} \to Y$ 是典范态射.

(i). 若 $f$ 满射，对 $Y^{\mathrm{an}}$ 的每个点 $y$，$f^{-1}(\psi(y))$ 是 $X$ 的非空闭部分；它含至少一个闭点，这证明 $f^{\mathrm{an}}$ 满射. 反之，若 $f^{\mathrm{an}}$ 满射，$f(X)$ 是 $Y$ 的局部\emph{可构部分} (EGA IV 1.8.4)，含 $Y$ 的所有闭点；故有 $f(X) = Y$.

(ii) 由 \ref{prop:XII.2.2} 得出.

(iii). 若 $f$ 是闭浸入，由 \ref{thm:XII.1.1} a) $f^{\mathrm{an}}$ 也如此. 反之，若 $f^{\mathrm{an}}$ 是闭浸入，由 \ref{prop:XII.3.1} (x) 和 \ref{prop:XII.3.2} (v) $f$ 也如此，因为说 $f$ 是闭浸入等价于说 $f$ 是紧合单态射 (EGA IV 8.11.5).

(iv). 显然 $f$ 是浸入时 $f^{\mathrm{an}}$ 也是. 反之假设 $f^{\mathrm{an}}$ 是浸入，设 $T$ 是 $X$ 在 $Y$ 中的像，$\overline{T}$ 是 $f$ 的概形闭包. 有 $f$ 的分解
\[
X \xrightarrow{i} \overline{T} \xrightarrow{j} Y,
\]
其中 $j$ 是闭浸入，$i$ 是典范态射，由此得到 $f^{\mathrm{an}}$ 的分解
\[
X^{\mathrm{an}} \xrightarrow{i^{\mathrm{an}}} \overline{T}^{\mathrm{an}} \xrightarrow{j^{\mathrm{an}}} Y^{\mathrm{an}}.
\]
由于 $T = f(X)$ 是 $Y$ 的局部可构部分 (EGA IV 1.8.4)，由 \ref{prop:XII.2.2} 有 $\overline{T}^{\mathrm{an}} = \overline{f^{\mathrm{an}}(X^{\mathrm{an}})}$. 由此得出 $i^{\mathrm{an}}(X^{\mathrm{an}})$ 是 $\overline{T}^{\mathrm{an}}$ 的开集，故 $i(X)$ 是 $\overline{T}$ 的开集. 考虑 $i$ 的典范分解
\[
X \xrightarrow{i_1} i(X) \xrightarrow{i_2} \overline{T}.
\]
态射 $i_1^{\mathrm{an}}$ 是紧合单态射，故由 \ref{prop:XII.3.2} (v) 和 \ref{prop:XII.3.1} (x) $i_1$ 也如此；这证明 $i_1$ 从而 $f$ 是浸入.

(v). 假设 $f$ 紧合，证明 $f^{\mathrm{an}}$ 也紧合. $f^{\mathrm{an}}$ 紧合的事实局部于 $Y^{\mathrm{an}}$，可假设 $Y$ 仿射. 由 Chow 引理 (EGA II 5.6.1)，可以找到 $Y$-射影概形 $X'$ 和射影满射态射
\[
g \colon X' \longrightarrow X.
\]
态射 $(fg)^{\mathrm{an}} = f^{\mathrm{an}} g^{\mathrm{an}}$ 射影故紧合，$g^{\mathrm{an}}$ 满射，由 \cite[第 1 章 §10]{exp12-1} 得出 $f^{\mathrm{an}}$ 紧合.

反之假设 $f^{\mathrm{an}}$ 紧合，证明 $f$ 也紧合. 由 \ref{prop:XII.3.1} (viii) $f$ 分离. 剩下证明 $f$ 万有闭，甚至只需证明 $f$ 闭；确实，对每个局部有限型 $Y$-概形 $Y'$，态射
\[
f_{(Y')} = h \colon X \times_Y Y' \longrightarrow Y'
\]
也闭，因为 $h^{\mathrm{an}}$ 紧合. 设 $T$ 是 $X$ 的闭部分；$f(T)$ 是局部可构集，且有
\[
f^{\mathrm{an}}(\phi^{-1}(T)) = \psi^{-1}(f(T)).
\]
由于 $f^{\mathrm{an}}$ 紧合，$\psi^{-1}(f(T))$ 是 $Y^{\mathrm{an}}$ 的闭部分，由 \ref{prop:XII.2.2} 得出
\[
\overline{\psi^{-1}(f(T))} = \psi^{-1}(\overline{f(T)}).
\]
这蕴含 $f(T) = \overline{f(T)}$，i.e. $f$ 闭故紧合.

(vi). 说态射有限与说它纤维有限紧合相同 (EGA III 4.4.2 和 \cite[第 19 号 §5]{exp12-4}). 由于 $f$ 的纤维有限的点集局部可构 (EGA IV 9.7.9)，$f$ 的诸纤维有限当且仅当 $f^{\mathrm{an}}$ 的诸纤维如此；故 (vi) 由 (v) 得出.

\begin{remark}\label{rem:XII.3.3}
a) 设 $f \colon X \to Y$ 是 $\mathbb{C}$ 上局部有限型概形的态射. $f^{\mathrm{an}}$ 是局部同构的事实不蕴含 $f$ 也如此. 确实，若 $f$ 平展，$f^{\mathrm{an}}$ 平展故是局部同构 \cite[第 13 号 §1]{exp12-4}，但 $f$ 未必如此.

b) 不假设 $f$ 有限型时断言 \ref{prop:XII.3.2} 不成立. 例如证明 $f^{\mathrm{an}}$ 可以是闭浸入而 $f$ 不是. 只需取 $X$ 为 $\mathbb{Z}$ 个 $\mathrm{Spec} \mathbb{C}$ 的和，$Y$ 为仿射直线，$f$ 为把 $X$ 的点送到 $Y$ 的构成离散部分的互异点的态射.
\end{remark}

\section{上同调比较定理与存在定理}

本号的对象是重新证明 \cite[第 2 号 定理 5 和 定理 6]{exp12-3} 的结果；后者把 $X$ 射影时在 \cite[第 12 号]{exp12-10} 中建立的定理推广到紧合概形的情形，并推广到相对情形. 关于解析空间上的相对紧合概形的更一般结果在 \cite[第八章 第 3 号]{exp12-7} 中证明.

回忆 \cite[第 12 号]{exp12-10} 中用的 Čech 上同调在代数情形与解析情形都等于通常上同调 (EGA III 1.4.1 和 \cite[II 5.10]{exp12-5}).

\refstepcounter{theorem}\noindent\textbf{\thesection.\arabic{theorem}.}\quad 设 $f \colon X \to Y$ 是 $\mathbb{C}$ 上局部有限型概形的态射，考虑交换图式
\[
\begin{tikzcd}[row sep=large, column sep=large]
X^{\mathrm{an}} \arrow[r,"{\phi}"] \arrow[d,"{f^{\mathrm{an}}}"] & X \arrow[d,"{f}"] \\
Y^{\mathrm{an}} \arrow[r,"{\psi}"] & Y.
\end{tikzcd}
\]
若 $F$ 是 $\mathcal{O}_X$-模，对每个整数 $p \geqslant 0$ 有态射
\[
R^p f_* F \xrightarrow{i} R^p f_*(\phi_* F^{\mathrm{an}}) \xrightarrow{j} R^p (f \cdot \phi)_* F^{\mathrm{an}} \xrightarrow{k} \psi_* (R^p f_*^{\mathrm{an}} F^{\mathrm{an}}),
\]
其中 $i$ 由典范态射 $F \to \phi_* F^{\mathrm{an}}$ 得出，$j$，$k$ 是 Leray 谱序列的"边缘同态". 与复合 $k \cdot j \cdot i$ 关联典范态射
\[
(4.1.1) \quad \theta_p \colon (R^p f_* F)^{\mathrm{an}} \longrightarrow R^p f_*^{\mathrm{an}} (F^{\mathrm{an}})
\]

\begin{theorem}\label{thm:XII.4.2}
设 $f \colon X \to Y$ 是 $\mathbb{C}$ 上局部有限型概形的紧合态射，$F$ 是凝聚 $\mathcal{O}_X$-模. 则对每个整数 $p \geqslant 0$，态射 (4.1.1)
\[
\theta_p \colon (R^p f_* F)^{\mathrm{an}} \longrightarrow R^p f_*^{\mathrm{an}} (F^{\mathrm{an}})
\]
是同构.
\end{theorem}

1) $f$ 射影的情形. 证明与 \cite[第 13 号]{exp12-10} 类似. 简要回忆. 归结到 $X$ 是 $Y$ 上方射影空间型 $\mathbf{P}^r_Y$ 的情形. 设 $\mathcal{Y} = Y^{\mathrm{an}}$，$\mathcal{P} = \mathbf{P}^r_\mathcal{Y}$；先证明
\[
f_*^{\mathrm{an}} \mathcal{O}_{\mathcal{P}} = \mathcal{O}_{\mathcal{Y}}, \quad R^p f_*^{\mathrm{an}}(\mathcal{O}_{\mathcal{P}}) = 0 \quad \text{对 } p > 0.
\]
为验证前面关系，可归结到 $\mathcal{Y}$ 是仿射空间 $E^n$ 的球 $B$ 的情形. 考虑 $\mathcal{P}$ 的由 $r + 1$ 个同构于 $B \times E^r$ 的开集组成的"标准覆盖" $\{U_i\}$. 由于这些开集是 Stein 的，对每个整数 $p \geqslant 0$ 有同构
\[
H^p(\{U_i\}, \mathcal{O}_{\mathcal{P}}) \xrightarrow{\sim} H^p(\mathcal{P}, \mathcal{O}_{\mathcal{P}})
\]
于是可以把结构层 $\mathcal{O}_{\mathcal{P}}$ 在开集 $U_i$ 及其交上的截面用 Laurent 级数表达；容易的计算证明
\[
H^0(\mathcal{P}, \mathcal{O}_{\mathcal{P}}) \xrightarrow{\sim} H^0(\mathcal{Y}, \mathcal{O}_{\mathcal{Y}}), \quad H^p(\mathcal{P}, \mathcal{O}_{\mathcal{P}}) = 0 \quad \text{对 } p > 0.
\]
证明于是照抄 \cite[第 12 号 引理 5]{exp12-10} 完成，上同调群换成上同调层.

2) $f$ 紧合的情形. 用 EGA III 3.1.2 归结到射影情形. 设 $K$ 是使 $\theta_p$ 对每个 $p \geqslant 0$ 是同构的凝聚 $\mathcal{O}_X$-模范畴. 只需证明：对每个正合序列 $0 \to F' \to F \to F'' \to 0$ (其中两项属于 $K$)，第三项也属于 $K$，$K$ 的对象的直和因子属于 $K$，且对 $X$ 的每个点 $x$，可以找到 $K$ 的对象 $F$ 使 $F_x \neq 0$.

第一个条件由五引理应用于下列行正合的交换图式得出：
\[
\begin{tikzcd}[row sep=large, column sep=large]
(R^p f_* F')^{\mathrm{an}} \arrow[r] \arrow[d] & (R^p f_* F)^{\mathrm{an}} \arrow[r] \arrow[d] & (R^p f_* F'')^{\mathrm{an}} \arrow[r] \arrow[d] & (R^{p+1} f_* F')^{\mathrm{an}} \\
R^p f_*^{\mathrm{an}} F'^{\mathrm{an}} \arrow[r] & R^p f_*^{\mathrm{an}} F^{\mathrm{an}} \arrow[r] & R^p f_*^{\mathrm{an}} F''^{\mathrm{an}} \arrow[r] & R^{p+1} f_*^{\mathrm{an}} F'^{\mathrm{an}},
\end{tikzcd}
\]
类似验证第二个条件.

为验证第三个条件，可限于 $X$ 是不可约概形、生成点为 $x$ 的情形. 可从一开始假设 $Y$ 诺特. 由 Chow 引理 (EGA II 5.6.1)，可以找到 $Y$-射影概形 $X'$ 和射影满射态射 $g \colon X' \to X$. 另一方面存在整数 $n$，使对每个 $p > 0$ 有 $R^p g_*(\mathcal{O}_{X'}(n)) = 0$，且典范态射 $g^{*} g_*(\mathcal{O}_{X'}(n)) \to \mathcal{O}_{X'}(n)$ 满射 (EGA III 2.2.1). 令 $F = g_*(\mathcal{O}_{X'}(n))$，层 $F$ 回答这个问题. 确实有 $F_x \neq 0$；而且 Leray 谱序列
\[
R^p f_*(R^q g_*(\mathcal{O}_{X'}(n))) \Longrightarrow R^{p+q} (f \cdot g)_* (\mathcal{O}_{X'}(n))
\]
退化，有同构
\[
R^p f_* F \xrightarrow{\sim} R^p (f \cdot g)_* (\mathcal{O}_{X'}(n)).
\]
如同代数情形有典范同构
\[
R^p f_*^{\mathrm{an}} F^{\mathrm{an}} \xrightarrow{\sim} R^p (f \cdot g)^{\mathrm{an}}_* (\mathcal{O}_{X'}(n)^{\mathrm{an}}),
\]
且图式
\[
\begin{tikzcd}[row sep=large, column sep=large]
(R^p f_* F)^{\mathrm{an}} \arrow[r,"{\sim}"] \arrow[d,"{\theta_p}"] & (R^p (f \cdot g)_* (\mathcal{O}_{X'}(n)))^{\mathrm{an}} \arrow[d,"{\psi_p}"] \\
R^p f_*^{\mathrm{an}} F^{\mathrm{an}} \arrow[r,"{\sim}"] & R^p (f \cdot g)^{\mathrm{an}}_* (\mathcal{O}_{X'}(n)^{\mathrm{an}})
\end{tikzcd}
\]
交换. 由 1) $\psi_p$ 是同构；故 $\theta_p$ 也同构，这完成证明.

\begin{corollary}\label{cor:XII.4.3}
设 $X$ 是 $\mathbb{C}$ 上紧合概形，$F$ 是凝聚 $\mathcal{O}_X$-模. 则对每个整数 $p \geqslant 0$，典范态射
\[
H^p(X, F) \longrightarrow H^p(X^{\mathrm{an}}, F^{\mathrm{an}})
\]
是同构.
\end{corollary}

\begin{theorem}\label{thm:XII.4.4}
设 $X$ 是 $\mathbb{C}$ 上紧合概形. 把凝聚 $\mathcal{O}_X$-模 $F$ 配给它在 $X^{\mathrm{an}}$ 上的逆像 $F^{\mathrm{an}}$ 的函子是范畴等价.
\end{theorem}

1) 该函子忠实全. 设 $F$ 和 $G$ 是两个凝聚 $\mathcal{O}_X$-模. 典范态射
\[
\mathrm{Hom}_{\mathcal{O}_X}(F, G) \longrightarrow \mathrm{Hom}_{\mathcal{O}_{X^{\mathrm{an}}}}(F^{\mathrm{an}}, G^{\mathrm{an}})
\]
等同于典范态射
\[
H^0(X, \mathcal{H}om_{\mathcal{O}_X}(F, G)) \longrightarrow H^0(X^{\mathrm{an}}, \mathcal{H}om_{\mathcal{O}_X}(F, G))
\]
(EGA III$_0$ 6.7.6). 由于 $\mathcal{H}om_{\mathcal{O}_X}(F, G)$ 凝聚，由 \ref{cor:XII.4.3} 该态射双射.

2) 该函子本质满射. 当 $X$ 射影时断言由 \cite[第 12 号 定理 3]{exp12-10} 得出. 一般情形用 Chow 引理 (EGA II 5.6.1) 归结到前述. 设 $X'$ 是 $\mathbb{C}$-射影概形，$f \colon X' \to X$ 是射影满射态射，$U$ 是 $X$ 的稠密开集，使 $f$ 诱导同构 $f^{-1}(U) \simeq U$. 对 $X$ 作诺特递归；可假设：对 $X^{\mathrm{an}}$ 上每个可以找到 $X$ 的与 $X$ 不同的闭部分 $Y$ (满足 $Y^{\mathrm{an}} \supset \mathrm{Supp} \mathcal{G}$) 的凝聚层 $\mathcal{G}$，存在 $X$ 上的凝聚层 $G$ 使有同构 $G^{\mathrm{an}} \simeq \mathcal{G}$.

设 $F$ 是 $\mathcal{O}_{X^{\mathrm{an}}}$ 上的凝聚模层，$K$ 和 $L$ 是由序列
\[
0 \longrightarrow K \longrightarrow F \longrightarrow f_*^{\mathrm{an}} f^{\mathrm{an}*} F \longrightarrow L \longrightarrow 0
\]
正合定义的凝聚层. 由于 $X'$ 射影，存在凝聚 $\mathcal{O}_{X'}$-模 $F'$ 使 $F'^{\mathrm{an}} \simeq f^{\mathrm{an}*} F$；由 \ref{thm:XII.4.2} 得出同构 $(f_* F')^{\mathrm{an}} \simeq f_*^{\mathrm{an}} f^{\mathrm{an}*} F$. 由于 $K|_{U^{\mathrm{an}}}$ 和 $L|_{U^{\mathrm{an}}}$ 为零，存在凝聚 $\mathcal{O}_X$-模 $K$ 和 $L$ 使有同构 $K^{\mathrm{an}} \simeq K$，$L^{\mathrm{an}} \simeq L$. 由 1) 态射 $f_*^{\mathrm{an}} f^{\mathrm{an}*} F \to L$ 来自唯一态射 $f_* F' \to L$；设 $I = \mathrm{Ker}(f_* F' \to L)$. 层 $F$ 是 $I^{\mathrm{an}}$ 被 $K^{\mathrm{an}}$ 的扩张，只需看这个扩张来自 $I$ 被 $K$ 的扩张的逆像. 故只需证明典范态射
\[
(\ast) \quad \mathrm{Ext}^q_{\mathcal{O}_X}(I, K)^{\mathrm{an}} \longrightarrow \mathrm{Ext}^q_{\mathcal{O}_{X^{\mathrm{an}}}}(I^{\mathrm{an}}, K^{\mathrm{an}}) \quad q \neq 1
\]
双射. 有同构 $\mathrm{Ext}^q_{\mathcal{O}_X}(I, K)^{\mathrm{an}} \xrightarrow{\sim} \mathrm{Ext}^q_{\mathcal{O}_{X^{\mathrm{an}}}}(I^{\mathrm{an}}, K^{\mathrm{an}})$ 对每个整数 $q \geqslant 0$ 成立 (EGA III$_0$ 12.3.5)，以及谱序列的态射
\[
\begin{tikzcd}[row sep=large, column sep=large]
H^p(X, \mathcal{E}xt^q_{\mathcal{O}_X}(I, K)) \arrow[r, Rightarrow] \arrow[d] & \mathrm{Ext}^{p+q}_{\mathcal{O}_X}(I, K) \arrow[d] \\
H^p(X^{\mathrm{an}}, \mathcal{E}xt^q_{\mathcal{O}_{X^{\mathrm{an}}}}(I^{\mathrm{an}}, K^{\mathrm{an}})) \arrow[r, Rightarrow] & \mathrm{Ext}^{p+q}_{\mathcal{O}_{X^{\mathrm{an}}}}(I^{\mathrm{an}}, K^{\mathrm{an}}).
\end{tikzcd}
\]
该态射是同构，因为由 \ref{cor:XII.4.3} 在 $E^{pq}_2$ 项上如此，这证明 $(\ast)$ 双射.

\begin{corollary}\label{cor:XII.4.5}
把紧合 $\mathbb{C}$-概形 $X$ 配给 $X^{\mathrm{an}}$ 的函子忠实全.
\end{corollary}

需证明：若 $X$ 和 $Y$ 是两个紧合 $\mathbb{C}$-概形，典范映射
\[
\mathrm{Hom}_{\mathbb{C}}(X, Y) \longrightarrow \mathrm{Hom}(X^{\mathrm{an}}, Y^{\mathrm{an}})
\]
双射. 给出 $X$ 到 $Y$ 的态射 (resp. $X^{\mathrm{an}}$ 到 $Y^{\mathrm{an}}$ 的态射) 等价于给出它的图，i.e. $X \times Y$ 的闭子概形 $Z$ (resp. $X^{\mathrm{an}} \times Y^{\mathrm{an}}$ 的闭解析子空间 $Z$)，使第一个投影 $X \times Y \to X$ (resp. $X^{\mathrm{an}} \times Y^{\mathrm{an}} \to X^{\mathrm{an}}$) 在 $Z$ 上的限制是同构. 由于给出 $X \times Y$ 的闭子概形 (resp. $X^{\mathrm{an}} \times Y^{\mathrm{an}}$ 的闭解析子空间) 等价于给出 $\mathcal{O}_{X \times Y}$ (resp. $\mathcal{O}_{X^{\mathrm{an}} \times Y^{\mathrm{an}}}$) 上的凝聚理想层，该推论由 \ref{thm:XII.4.4} 得出.

\begin{corollary}\label{cor:XII.4.6}
设 $X$ 是 $\mathbb{C}$ 上紧合概形. 把 $X$ 上方的有限概形 (resp. 有限平展概形) $X'$ 配给 $X'^{\mathrm{an}}$ 的函子，是 $X$ 上方有限概形 (resp. 有限平展概形) 范畴到 $X^{\mathrm{an}}$ 上方有限 (resp. 有限平展) 解析空间范畴的等价.
\end{corollary}

确实，给出有限态射 $X' \to X$ (resp. $X'^{\mathrm{an}} \to X^{\mathrm{an}}$) 等价于给出 $\mathcal{O}_X$ (resp. $\mathcal{O}_{X^{\mathrm{an}}}$) 上的凝聚代数层 \cite[第 19 号 §5 定理 2]{exp12-4}. 该推论在非 resp. 情形由 \ref{thm:XII.4.4} 得出，resp. 情形考虑到 \ref{prop:XII.3.1} (iii) 由此推出.

\section{平展覆盖的比较定理}

\refstepcounter{theorem}\addtocounter{theorem}{-1}\noindent\textbf{\thesection.0.}\quad 明确\emph{解析空间的有限覆盖}概念. 若 $X$ 是解析空间，称 $X$ 上方有限的解析空间 $X'$ 为 $X$ 的有限覆盖\index[term]{jiexikongjiandeyouxianfugai@解析空间的有限覆盖}，若 $X'$ 的每个不可约分量支配 $X$ 的某个不可约分量.

\begin{theorem}\label{thm:XII.5.1}
[Riemann 存在定理]\index[term]{Riemann cunzaidingli@Riemann 存在定理} 设 $X$ 是 $\mathbb{C}$ 上局部有限型概形，$X^{\mathrm{an}}$ 是 $X$ 的关联解析空间. 把 $X$ 的每个有限平展覆盖 $X'$ 配给 $X'^{\mathrm{an}}$ 的函子 $\Psi$，是 $X$ 的有限平展覆盖范畴到 $X^{\mathrm{an}}$ 的有限平展覆盖范畴的等价.
\end{theorem}

1) 函子 $\Psi$ 忠实全. 设 $X'$ 和 $X''$ 是 $X$ 的两个有限平展覆盖，证明典范映射
\[
(\ast) \quad \mathrm{Hom}_X(X', X'') \longrightarrow \mathrm{Hom}_{X^{\mathrm{an}}}(X'^{\mathrm{an}}, X''^{\mathrm{an}})
\]
双射. 可假设 $X'$ 连通. 给出 $X'$ 到 $X''$ 的 $X$-态射等价于给出 $X' \times_X X''$ 的连通分量 $X_i$，使第一个投影诱导的态射 $X_i \to X'$ 是同构. 由于 $X' \times_X X''$ 的连通分量与 $X'^{\mathrm{an}} \times_{X^{\mathrm{an}}} X''^{\mathrm{an}}$ 的连通分量一一对应 (\cref{cor:XII.2.6})，且态射 $X_i \to X'$ 是同构当且仅当 $X_i^{\mathrm{an}} \to X'^{\mathrm{an}}$ 如此，这证明 $(\ast)$ 双射.

2) 函子 $\Psi$ 本质满射. 设 $X'$ 是 $X^{\mathrm{an}}$ 的有限平展覆盖，证明存在 $X$ 的平展覆盖 $X'$，使有同构 $X'^{\mathrm{an}} \xrightarrow{\sim} X'$. 考虑到 1) 问题局部于 $X$，可假设 $X$ 仿射.

a) 归结到 $X$ 正规的情形. 可假设 $X$ 既约. 假设定理对 $X_{\mathrm{r\'ed}}$ 已证. 把 $X$ 的有限平展覆盖 $X'$ 配给 $X_{\mathrm{r\'ed}}$ 的有限平展覆盖 $X_{\mathrm{r\'ed}}'^{\mathrm{an}}$ 的函子此时是等价. 由于它由 $\Psi$ 与把 $X^{\mathrm{an}}$ 的有限平展覆盖配给它在 $X_{\mathrm{r\'ed}}^{\mathrm{an}}$ 上的逆像的函子 $\Theta$ 复合得到，且 $\Theta$ 忠实全，这证明 $\Psi$ 是范畴等价.

可假设 $X$ 正规. 设 $\widetilde{X}$ 是 $X$ 的正规化，$p \colon \widetilde{X} \to X$ 是典范态射. 由于 $p$ 有限，$p$ 对平展覆盖范畴是有效下降态射 (\cref{thm:IX.4.7}). 假设定理对 $\widetilde{X}$ 已证，若令 $\widetilde{X}' = X' \times_{X^{\mathrm{an}}} \widetilde{X}^{\mathrm{an}}$，存在 $\widetilde{X}$ 的平展覆盖 $\widetilde{X}'$ 和同构 $\widetilde{X}'^{\mathrm{an}} \simeq \widetilde{X}'$. 由 1) 得出 $\widetilde{X}'$ 上的自然下降数据提升为相对 $\widetilde{X} \to X$ 的 $\widetilde{X}'$ 上的下降数据；这证明存在 $X$ 的平展覆盖 $X'$，使有同构 $i \colon X'^{\mathrm{an}} \times_{X^{\mathrm{an}}} \widetilde{X}^{\mathrm{an}} \simeq \widetilde{X}'$，它在 $\widetilde{X}^{\mathrm{an}} \times_{X^{\mathrm{an}}} \widetilde{X}^{\mathrm{an}}$ 的两个投影下的逆像相同. 由 IX 3.2 (其证明在解析情形有效)，态射 $\widetilde{X}^{\mathrm{an}} \to X^{\mathrm{an}}$ 对平展覆盖范畴是下降态射，从而 $i$ 来自同构 $X'^{\mathrm{an}} \simeq X'$.

b) 归结到 $X$ 正则的情形. 设 $U$ 是 $X$ 的正则点的开集，$i \colon U \to X$，$i^{\mathrm{an}} \colon U^{\mathrm{an}} \to X^{\mathrm{an}}$ 是典范态射；由于 $X$ 正规，有 $\mathrm{codim}(X - U, X) \geqslant 2$. 假设存在 $U$ 的平展覆盖 $U'$ 使 $U'^{\mathrm{an}} \simeq X'|_{U^{\mathrm{an}}}$，证明此时 $U'$ 延拓成 $X$ 的平展覆盖 $X'$，使 $X'^{\mathrm{an}} \simeq X'$. 只需看 $U'$ 延拓成 $X$ 的平展覆盖 $X'$；确实将有同构 $X'^{\mathrm{an}}|_{U^{\mathrm{an}}} \simeq X'|_{U^{\mathrm{an}}}$；但若 $F$ 和 $G$ 是 $\mathcal{O}_{X^{\mathrm{an}}}$ 上分别定义 $X'$ 和 $X'^{\mathrm{an}}$ 的凝聚代数层，$X$ 正规且 $\mathrm{codim}(X - U, X) \geqslant 2$ 的事实蕴含典范态射
\[
F \longrightarrow i_*^{\mathrm{an}}(F|_{U^{\mathrm{an}}}), \quad G \longrightarrow i_*^{\mathrm{an}}(G|_{U^{\mathrm{an}}})
\]
是同构 \cite[第 3 号 命题 4]{exp12-11}. 由此得出 $F$ 和 $G$ 从而 $X'^{\mathrm{an}}$ 和 $X'$ 同构.

设 $\phi \colon X^{\mathrm{an}} \to X$ 是典范态射. 由于把 $U'$ 延拓到 $X$ 的问题局部于 $X$，只需证明：对 $X^{\mathrm{an}} - U^{\mathrm{an}}$ 的每个点 $y$，$U_{\phi(y)} = U \times_X \mathrm{Spec} \mathcal{O}_{X,\phi(y)}$ 的平展覆盖 $U'_{\phi(y)} = U' \times_X \mathrm{Spec} \mathcal{O}_{X,\phi(y)}$ 延拓到 $\mathrm{Spec} \mathcal{O}_{X,\phi(y)}$. 设 $H$ 是定义 $U'$ 的 $\mathcal{O}_U$-凝聚代数. 典范态射
\[
\alpha \colon (i_* H)^{\mathrm{an}} \longrightarrow i_*^{\mathrm{an}}(H^{\mathrm{an}}) = F
\]
定义 $\mathrm{Spec} \mathcal{O}_{X^{\mathrm{an}},y}$ 上模层的态射：
\[
\alpha_y \colon (i_* H)^{\mathrm{an}}_y \longrightarrow F_y,
\]
它在 $U_y = U_{\phi(y)} \times_{\mathrm{Spec} \mathcal{O}_{X,\phi(y)}} \mathrm{Spec} \mathcal{O}_{X^{\mathrm{an}},y}$ 上的限制是同构. 但这证明 $H|_{U_y}$ 平凡，故 $U'_{\phi(y)}$ 延拓到 $\mathrm{Spec} \mathcal{O}_{X,\phi(y)}$.

c) $X$ 仿射正则的情形. 设
\[
j \colon X \longrightarrow P
\]
是 $X$ 的紧化，其中 $P$ 是射影 $\mathbb{C}$-概形，$j$ 是支配开浸入. 由奇点消解定理\index[term]{qidianxiaojie@奇点消解} \cite{exp12-8}，可以找到正则概形 $R$、射影态射 $r \colon R \to P$，使 $r$ 诱导同构 $r^{-1}(X) \simeq X$，且 $r^{-1}(X)$ 是 $R$ 中正规交叉除子的余集. 设
\[
k \colon X \longrightarrow R
\]
是典范浸入. 将证明存在 $R^{\mathrm{an}}$ 的延拓平展覆盖 $X'^{\mathrm{an}}$ 的有限正规覆盖 (5.0) $R'$. 由下面\cref{prop:XII.5.3}，这样的覆盖唯一；延拓 $X'^{\mathrm{an}}$ 的问题于是局部于 $R^{\mathrm{an}}$ 在 $R^{\mathrm{an}} - X^{\mathrm{an}}$ 附近的邻域. 而 $R^{\mathrm{an}} - X^{\mathrm{an}}$ 的每个点有同构于仿射空间 $E^n$ 的球的开邻域 $V$，使 $V - V \cap X^{\mathrm{an}}$ 由前 $p$ 个坐标函数 $z_1, \ldots, z_p$ (其中 $0 \leqslant p \leqslant n$) 的消没定义. $U = V \cap X^{\mathrm{an}}$ 的基本群同构于 $\mathbb{Z}^p$，$U$ 的每个平展覆盖是形如
\[
U'' = U[T_1, \ldots, T_p]/(T_1^{n_1} - z_1, \ldots, T_p^{n_p} - z_p)
\]
的覆盖除以 $U''$ 的 伽罗瓦群 $\mathbb{Z}/n_1\mathbb{Z} \times \cdots \times \mathbb{Z}/n_p\mathbb{Z}$ 的子群 $H$ 的商，其中诸 $n_i$ 是 $> 0$ 的整数. 而 $U''$ 延拓成正则覆盖
\[
V'' = V[T_1, \ldots, T_p]/(T_1^{n_1} - z_1, \ldots, T_p^{n_p} - z_p),
\]
$H$ 作用在其上，$V''$ 除以 $H$ 的商是所求延拓.

证明于是由 \ref{cor:XII.4.6} 完成. 覆盖 $R'$ 来自 $R$ 的有限覆盖 $R'$；$R'$ 在 $X$ 上的限制是 $X$ 的覆盖 $X'$，使 $X'^{\mathrm{an}} \simeq X'$，且由 \ref{prop:XII.3.1} (iii) $X'$ 是 $X$ 的平展覆盖.

\begin{corollary}\label{cor:XII.5.2}
设 $X$ 是 $\mathbb{C}$ 上连通局部有限型概形，$\phi \colon X^{\mathrm{an}} \to X$ 是典范态射，$x$ 是 $X^{\mathrm{an}}$ 的点. 设 $\pi_1(X^{\mathrm{an}}, x)$ 是拓扑空间 $X^{\mathrm{an}}$ 在点 $x$ 的基本群，$\pi_1(X, \phi(x))$ 是概形 $X$ 在点 $\phi(x)$ 的基本群 (V 7). 则 $\pi_1(X, \phi(x))$ 典范同构于 $\pi_1(X^{\mathrm{an}}, x)$ 关于有限指数子群拓扑的完备化.
\end{corollary}

设 $\mathcal{C}$ 是 $X^{\mathrm{an}}$ 的有限平展覆盖范畴，$F$ 是从 $\mathcal{C}$ 到 $\mathrm{Ens}$ 的、把 $X^{\mathrm{an}}$ 的每个有限平展覆盖 $X'$ 配给 $X'$ 的位于 $x$ 上方的点集的函子，设 $\widehat{\pi_1}(X^{\mathrm{an}}, x)$ 是按 V 4 所述与 $\mathcal{C}$ 和 $F$ 关联的 pro-有限群. 由于 $X^{\mathrm{an}}$ 的每个有限平展覆盖是万有覆盖除以有限指数子群的商，$\widehat{\pi_1}(X^{\mathrm{an}}, x)$ 正是 $\pi_1(X^{\mathrm{an}}, x)$ 关于有限指数子群拓扑的完备化. 该推论于是由 \ref{thm:XII.5.1} 和 V 6.10 得出.

\begin{proposition}\label{prop:XII.5.3}
设 $X$ 是正规解析空间，$Y$ 是 $U = X - Y$ 在 $X$ 中稠密的闭解析子集. 则把 $X$ 的每个有限正规覆盖 (5.0) $X'$ 配给它在 $U$ 上的限制的函子忠实全.
\end{proposition}

设 $X'$ 和 $X''$ 是 $X$ 的两个有限正规覆盖. 需证明典范映射
\[
\mathrm{Hom}_X(X', X'') \longrightarrow \mathrm{Hom}_U(X'|_U, X''|_U)
\]
双射. 设 $u$，$v$ 是 $X'$ 到 $X''$ 的两个在 $U$ 上的限制相同的 $X$-态射，证明 $u = v$. 态射 $u$ 和 $v$ 在稠密开集 $U \times_X X'$ 上重合，故在底拓扑空间上重合. 由 \cite[第 19 号 §4 推论 5]{exp12-4} 这证明 $u = v$.

现在设 $u$ 是 $X'|_U$ 到 $X''|_U$ 的 $U$-态射，证明它延拓到整个 $X'$. 可假设 $X'$ 正则. 确实，$X'$ 正规，可以找到余集为余维数 $\geqslant 2$ 的解析部分的 $X$ 的开集 $V$，使 $X' \times_X V = V'$ 正则. 设 $V'' = X'' \times_X V$，假设命题对 $V$ 已证. 考虑交换图式
\[
\begin{tikzcd}[row sep=large, column sep=large]
V' \arrow[r,"{i'}"] & X' \arrow[d,"{f'}"] \\
V'' \arrow[r,"{i''}"'] & X'' \arrow[d,"{f''}"] \\
V \arrow[r,"{i}"'] & X
\end{tikzcd}
\]
与 $u$ 关联 $\mathcal{O}_V$-代数的态射 $g_*'' \mathcal{O}_{V''} \to g_*' \mathcal{O}_{V'}$，由此得到态射
\[
i_* g_*'' \mathcal{O}_{V''} \longrightarrow i_* g_*' \mathcal{O}_{V'}.
\]
考虑到同构 $i'_* \mathcal{O}_{V'} \simeq \mathcal{O}_{X'}$，$i''_* \mathcal{O}_{V''} \simeq \mathcal{O}_{X''}$ \cite[第 3 号 命题 4]{exp12-11}，由此得到 $\mathcal{O}_X$-代数的态射
\[
f_*'' \mathcal{O}_{X''} \longrightarrow f_*' \mathcal{O}_{X'},
\]
从而得到所求态射 $X' \to X''$.

此后假设 $X'$ 正则. 设 $U' = U \times_X X'$，$Y' = X' - U'$. 把 $Y'$ 看作 $X'$ 的既约解析子空间；若 $Y'_1$ 是 $Y'$ 的奇异闭集，有 $\dim Y'_1 < \dim Y'$ \cite[第 20 号 D. 定理 3]{exp12-4}. 对 $\dim Y'$ 作递归，可假设 $Y'$ 光滑. 由于只需把 $u$ 延拓到 $Y'$ 每个点的开邻域，由隐函数定理可假设 $X'$ 是仿射空间 $E^n$ 的球，$Y'$ 由前 $p$ 个坐标函数 $z_1, \ldots, z_p$ (其中 $0 \leqslant p \leqslant n$) 的消没定义的闭集.

把 $u$ 关联 $p \colon X' \times_X X'' \to X'$ 在 $U'$ 上方的截面 $s$. 可以限制 $X'$，使 $A = p_*(\mathcal{O}_{X' \times_X X''})$ 由 $\Gamma(X', A)$ 的元素 $x_1, \ldots, x_q$ 生成. 设 $u_1, \ldots, u_q \in \Gamma(U', \mathcal{O}_{X'})$ 是 $s$ 对 $x_1|_{U'}, \ldots, x_q|_{U'}$ 的像. 说 $s$ 延拓到 $X'$ 等价于说诸 $u_1, \ldots, u_q$ 延拓成 $\Gamma(X', \mathcal{O}_{X'})$ 的截面. 但 $f$ 有限，每个 $u_i$ 是 $z_1, \ldots, z_p$ 的 Laurent 级数 (系数是 $z_{p+1}, \ldots, z_n$ 的幂级数)，满足整依赖关系. 由此得出 $u_i$ 有界故是 $z_1, \ldots, z_n$ 的幂级数，从而延拓到 $X'$.

可以问 \ref{prop:XII.5.3} 引入的函子是否是范畴等价. 借助 Grauert-Remmert 定理 \cite{exp12-6} 可以回答这个问题，下面用奇点消解给出它的证明. 在 \cite{exp12-8} 之前也可以用 Grauert-Remmert 定理证明 \ref{thm:XII.5.1}；正是这样做的.

\begin{theorem}\label{thm:XII.5.4}
[Grauert-Remmert 定理]\index[term]{Grauert-Remmert dingli@Grauert-Remmert 定理} 设 $X$ 是正规解析空间，$Y$ 是 $U = X - Y$ 在 $X$ 中稠密的闭解析子集. 设 $U'$ 是 $U$ 的有限正规覆盖；假设存在 $X$ 的稀疏闭解析部分 $S$，使 $U'$ 在 $U - U \cap S$ 上的限制平展. 则存在 $X$ 的延拓 $U'$ 的有限正规覆盖 $X'$，且 $X'$ 在同构意义下唯一.
\end{theorem}

唯一性由 \ref{prop:XII.5.3} 得出. 延拓 $U'$ 的问题于是局部于 $X$. 可假设 $U$ 正则且 $U'$ 在 $U$ 上平展. 确实 $U$ 的正则点集是 $X$ 中稠密、余集为解析部分的开集 $V$ \cite[第 20 号 D 定理 2]{exp12-4}，只需把 $U$ 换成开集 $V - V \cap S$.

设 $y$ 是 $X - U$ 的点，证明可以把 $U'$ 延拓到 $y$ 的邻域. 把 $X$ 限制到 $y$ 的开邻域，由奇点消解定理 \cite{exp12-8} 可以找到正则解析空间 $X_1$、射影态射 $f \colon X_1 \to X$ (在 $U$ 上的限制诱导同构 $U_1 = f^{-1}(U) \simeq U$)，使 $U_1$ 是 $X_1$ 中正规交叉除子的余集. 证明 $U'$ 延拓成 $X_1$ 的有限正规覆盖. 由于问题局部于 $X_1$，可假设 $X_1$ 是仿射空间 $E^n$ 的球，$X_1 - U_1$ 由前 $p$ 个坐标函数 $z_1, \ldots, z_p$ (其中 $0 \leqslant p \leqslant n$) 的消没定义. $U_1$ 的平展覆盖 $U'$ 是形如
\[
U_2 = U_1[T_1, \ldots, T_p]/(T_1^{n_1} - z_1, \ldots, T_p^{n_p} - z_p)
\]
的覆盖除以 $U_2$ 的 伽罗瓦群的子群 $H$ 的商. 覆盖 $U_2$ 延拓成覆盖
\[
X_2 = X_1[T_1, \ldots, T_p]/(T_1^{n_1} - z_1, \ldots, T_p^{n_p} - z_p)
\]
$H$ 作用在其上，$X_2/H$ 把 $U'$ 延拓到 $X_1$.

记 $X'_1$ 为 $X_1$ 的延拓 $U'$ 的有限正规覆盖，$F_1$ 为 $F_1$ 定义的 $\mathcal{O}_{X_1}$-凝聚代数. 由 Grauert-Remmert 有限性定理 \cite[第 15 号 \cref{thm:XII.1.1}]{exp12-4}，$f_* F_1$ 是 $\mathcal{O}_X$-凝聚代数. 它对应 $X$ 的有限覆盖 $X'$，而且 $X'$ 正规，因为 $X'_1$ 正规，且 $X'$ 是所求 $U'$ 的延拓.

\begin{remark}\label{rem:XII.5.5}
在断言 \ref{thm:XII.5.4} 中，不能删除关于 $U' \to U$ 不平展的点集的条件. 例如取 $X$ 为复平面的单位圆盘，$U$ 为 $X$ 中原点的余集，$U' = U[T]/(T^2 - \sin 1/z)$，其中 $z$ 是 $X$ 上的坐标函数. 则 $U'$ 是 $U$ 的不延拓到 $X$ 的有限正规覆盖. 确实假设 $U'$ 延拓成 $X$ 的有限覆盖 $X'$；则 $X' \to X$ 不平展的点集是含所有满足 $\sin 1/z = 0$ 的点 $z$ 的闭解析集，这荒谬.

然而，当 $\mathrm{codim}(X - U, X) \geqslant 2$ 时可以删除关于 $U' \to U$ 的奇异点集的条件. 可假设 $U$ 正则. $U' \to U$ 不平展的点集是 $U$ 的除子，由 Remmert-Stein 定理 \cite[定理 3]{exp12-9} 得出它是 $X$ 的除子在 $U$ 上的迹. 此时，若 $A$ 是使 $U' = \mathrm{Spec}^{\mathrm{an}}(A)$ 的 $\mathcal{O}_U$-凝聚代数，$i \colon U \to X$ 是典范态射，只需取 $X' = \mathrm{Spec}^{\mathrm{an}}(i_* A)$；确实知道 $i_* A$ 凝聚 \cite[第 1 号 定理 1]{exp12-11}.

\begin{remark}[2003 年 (MR) 增补]
存在有限表现的非平凡群 $G$，没有与 $G$ 不同的有限指数子群，例如 Higman 群 (参见 J.-P. Serre，Arbres et amalgames，Ast\'erisque N\textdegree 46，命题 6，第一章 §1). 于是，若这样的群实现为 $\mathbb{C}$ 上概形 $V$ (说光滑且射影) 的拓扑基本群，拓扑空间 $V_{\mathrm{top}}$ (即 $V$ 的底拓扑空间) 不单连通，但 $V$ 代数单连通. 目前不知道这样的 $V$. 但指出 D. Toledo 构造了 $\mathbb{C}$ 上光滑射影概形 $V$，它的拓扑基本群不关于有限指数子群拓扑分离 (自然态射 $\pi_1(V_{\mathrm{top}}) \to \pi_1(V_{\mathrm{alg}})$ 不单射). [D. Toledo，Projective varieties with non-residually finite fundamental group，Publ. Math. IHES 77 (1993)，p. 103--119.]
\end{remark}
\end{remark}

\section{文献}\printbibliography[heading=none]


